\documentclass{jsarticle}
\usepackage{amsmath}
\usepackage{amssymb}
\begin{document}
\title{}
\author{}
\date{}
\maketitle
　Time has frozen in constant entropy of zeta function, things becoming also being frozen.
Past to go after time while things had frozen, while in there are things have frozen,
If relating past time go through while in there are information frozen,
Energy of things while time has frozen, go past times thing becoming constant,
Energy while in no work,

物事が止まっている間、　時が未来へ行っている間、
相対性理論が空間の先へ時が行くと、その間事物は止まっている。
物事は先へは行っていない。

まとめると、情報量の密度エネルギー$ m(x) $
$${d \over df}F(x)=m(x)$$
が、時が経つと拡散している。
時が先へ進むと、情報量密度エネルギーは拡散する。
すべては、タイムマシーンが、どのように機能するかである。時空がこの密度エネルギーが拡散するか動的に止まっているかで、時の流れを変えられる。すべては、ヒッグス場の方程式が、この密度エネルギーの大きさを変えるかで、時の流れを変えられて、相対性理論の仕組みで、光速とは関係なく、ニュートン方程式の剛体力学とも関係なく、時空をこのヒッグス場の方程式の密度エネルギーの大きさを変えるだけで、タイムマシーンは出来る。すべては、人体に存在しているリボ核酸の機能から、数学と物理学、医学の応用で、情報空間のビッグバーンは、生成される。
密度エネルギーの大きさが拡散すると時間は止まっている。密度エネルギーが止まっていると時間は流れる。タイムマシーンはこの密度エネルギーによって作られる。
時間はこの真逆の時の流れで作られている。時間が止まっていると、光速とは関係なく事物の移動が光速以上に先に進む。時間と空間が関係なくタイムマシーンの移動が出来る。すべての時間によるタイムマシーンは、ヒッグス場の密度エネルギーの大きさで作られる。
すべては、熱力学の第2法則と第1法則の情報量増大とエントロピー不変量が
ダークマターとヒッグス場の方程式が、この時の見かけの理論となっている。
相対性理論では時間は存在しているが、量子力学では時間は止まっている。
すべてはゼータ関数が、この相対性理論と量子力学とを統合する理論になっている。
アインシュタイン博士は、この静止宇宙を考えていた。
すべては、3次元多様体とAdS5時空のアーベル多様体が、このゼータ関数とヒッグス場の方程式を、4次元多様体とミンコフスキー時空を包み込んでいる。
ブラックホールとホワイトホールの関係が、宇宙と異次元空間の原子と遷移元素の光量子仮説の関係を作っている。





\end{document}
